% eksperymentalne tłumaczenie na włoski

\section{Pole powierzchni} % lang-pl
\section{Area} % lang-it
\label{sekcja_pole_powierzchni}

Pole powierzchni opisuje, jak duża jest dana figura płaska. % lang-pl
L'area descrive quanto è grande una data figura piana. % lang-it
Opiszemy, jak na przestrzeni wieków zmieniać się będzie matematyczny opis pola, skupiając się na rozcięciach. % lang-pl
Descriveremo come, nel corso dei secoli, cambierà la descrizione matematica dell'area, concentrandoci sulle dissezioni. % lang-it
Sekcję zwieńczy twierdzenie Wallace'a-Bolyaia-Gerwiena, charakteryzujące figury o równych polach. % lang-pl
La sezione si concluderà con il teorema di Wallace-Bolyai-Gerwien, che caratterizza le figure di area uguale. % lang-it
Miarę Lebesgue'a wzmiankujemy krótko na koniec w formie ciekawostki, nie będzie nam do niczego potrzebna. % lang-pl
Accenneremo brevemente alla misura di Lebesgue alla fine, come curiosità: non ci servirà a nulla. % lang-it

Euklides nawet nie zdefiniuje on figur o równych polach, a tym bardziej samego pola jako liczby. % lang-pl
Euclide non definirà nemmeno le figure di area uguale, tanto meno l'area stessa come numero. % lang-it
(Będzie unikać problemu kwadratury koła, ale w księdze XII napisze, że pola dwóch kół mają się do siebie tak, jak kwadraty ich promieni. % lang-pl
(Eviterà il problema della quadratura del cerchio, ma nel libro XII scriverà che le aree di due cerchi stanno tra loro come i quadrati dei loro raggi. % lang-it
Czytaj dalej w sekcji \ref{sekcja_przyblizenia_pi}.) % lang-pl
Continua a leggere nella sezione \ref{sekcja_przyblizenia_pi}.) % lang-it
Ale czytając jego dowody możemy domyślić się, że równopolowe figury mają następujące właściwości: % lang-pl
Ma leggendo le sue dimostrazioni possiamo intuire che le figure equivalenti (equal-area) hanno le seguenti proprietà: % lang-it
\begin{enumerate}
\item przystające figury mają równe pola; % lang-pl
\item le figure congruenti hanno aree uguali; % lang-it
\item suma, różnica i połówki figur o równych polach mają równe pola; % lang-pl
\item la somma, la differenza e le metà di figure di area uguale hanno aree uguali; % lang-it
\item całość jest większa od części i wreszcie % lang-pl
\item il tutto è maggiore della parte e, infine, % lang-it
\item boki kwadratów o równych polach są równej długości. % lang-pl
\item i lati di quadrati di area uguale hanno la stessa lunghezza. % lang-it
\end{enumerate}

Czwarty punkt jest wnioskiem z trzeciego. % lang-pl
Il quarto punto è una conseguenza del terzo. % lang-it
Hartshorne \cite[s. 196--204]{hartshorne_2000} podczas formalizacji pojęcia pola dla płaszczyzny Hilberta wspomni aksjomat de Zolta: % lang-pl
Hartshorne \cite[p. 196--204]{hartshorne_2000}, formalizzando il concetto di area per il piano di Hilbert, menzionerà l'assioma di De Zolt: % lang-it

\begin{axiom}
[Z, de Zolta] % lang-pl
[Z, di De Zolt] % lang-it
\index{aksjomat!de Zolta}% % lang-pl
\index{assioma!di De Zolt}% % lang-it
    Jeśli $Q$ jest figurą zawartą wewnątrz figury $P$ tak, że różnica $P \setminus Q$ ma niepuste wnętrze, to figury $P$ oraz $Q$ nie są równopolowe. % lang-pl
    Se $Q$ è una figura contenuta all'interno di una figura $P$ in modo tale che la differenza $P \setminus Q$ abbia interno non vuoto, allora le figure $P$ e $Q$ non sono equivalenti. % lang-it
\end{axiom}

Starożytni oczywiście nie będą znać teorii miary i postąpią inaczej: wykorzystają zapomniane przez ich potomków rozcięcia. % dissections. % lang-pl
Gli antichi, ovviamente, non conosceranno la teoria della misura e procederanno diversamente: utilizzeranno le dissezioni, dimenticate dai loro posteri. % lang-it
Technika ta pojawi się wielokrotnie, na przykład w dowodzie (I.35), że równoległoboki o~tej samej podstawie i równych wysokościach mają równe pola. % lang-pl
Questa tecnica comparirà più volte, per esempio nella dimostrazione (I.35) che i parallelogrammi con la stessa base e altezze uguali hanno aree uguali. % lang-it
W 1873 roku Henry Perigal (a w 1917 inny Henry, Dudeney) poda dowód twierdzenia Pitagorasa wykorzystujący rozcięcia. % lang-pl
Nel 1873 Henry Perigal (e nel 1917 un altro Henry, Dudeney) fornirà una dimostrazione del teorema di Pitagora basata sulle dissezioni. % lang-it
\index[persons]{Dudeney, Henry}%
\index[persons]{Perigal, Henry}%
Dokładnie opiszą to Eves \cite[s. 194-239]{eves1_1972} i (z nieco inną notacją) Hartshorne \cite[s. 212-221]{hartshorne_2000}. % lang-pl
Lo descriveranno in dettaglio Eves \cite[p. 194-239]{eves1_1972} e (con una notazione leggermente diversa) Hartshorne \cite[p. 212-221]{hartshorne_2000}. % lang-it

Niech $P, Q$ będą jakimiś figurami o tej samej powierzchni. % lang-pl
Siano $P, Q$ due figure di area uguale. % lang-it
Będziemy mówić, że są równopolowe (\emph{equivalent}) i pisać $P \simeq Q$. % lang-pl
Diremo che sono equivalenti (\emph{equivalent}) e scriveremo $P \simeq Q$. % lang-it
Oznaczenie $P \cong Q$ zachowujemy dla zwykłego przystawania. % lang-pl
Riserviamo la notazione $P \cong Q$ per la consueta congruenza. % lang-it

\begin{definition}
    \label{definicja_rownowazne_przez_dodawanie}%
    Dwa wielokąty $P$, $Q$ nazywamy równoważnymi przez rozcinanie, jeśli można rozciąć każdy na skończoną liczbę wielokątnych kawałków o rozłącznych wnętrzach:   % lang-pl
    Due poligoni $P$, $Q$ si dicono equiscomponibili se si può ritagliare ciascuno di essi in un numero finito di pezzi poligonali con interni disgiunti: % lang-it
    \begin{equation}
        P = \bigcup_{k=1}^n P_k, \quad\quad\quad Q = \bigcup_{k=1}^n Q_k,
    \end{equation}
    że kawałek $P_k$ przystaje do $Q_k$ (istnieje pewna izometria $f_k$ taka, że $f_k[P_k] = Q_k$). % lang-pl
    tali che il pezzo $P_k$ sia congruente a $Q_k$ (esiste cioè un'isometria $f_k$ tale che $f_k[P_k] = Q_k$). % lang-it
    Oznaczamy to $P \cong_+ Q$. % lang-pl
    Lo denotiamo con $P \cong_+ Q$. % lang-it
\end{definition}

Eves będzie mówić o przystawaniu przez dodawanie, ale w użyciu będą inne określenia, takie jak \emph{equidecomposable} albo \emph{scissors-congruent}. % lang-pl
Eves parlerà di congruenza per addizione, ma saranno in uso anche altri termini, come \emph{equidecomposable} oppure \emph{scissors-congruent}. % lang-it
Rozcinanie wielokątów będzie przewijać się przez wiele artykułów w Delcie: % lang-pl
Le dissezioni di poligoni ricorreranno in molti articoli su Delta: % lang-it
Joanny Jaszuńskiej w $\Delta_{15}^{8}$, % {https://www.deltami.edu.pl/2015/08/nozyczki-matematyczne/} % lang-pl
di Joanna Jaszuńska in $\Delta_{15}^{8}$, % lang-it
Marka Kordosa w $\Delta_{17}^4$, % {https://www.deltami.edu.pl/2017/04/dedukcja-lokalna-na-przykladzie/} % lang-pl
di Marek Kordos in $\Delta_{17}^4$, % lang-it
Bartłomieja Bzdęgi w $\Delta_{20}^6$ (rozcinanie na równoległoboki), % \todofoot{https://www.deltami.edu.pl/2020/06/rownoleglobok/ zad. 8} % lang-pl
di Bartłomiej Bzdęga in $\Delta_{20}^6$ (dissezione in parallelogrammi), % lang-it
Miłosza Kwiatkowskiego w $\Delta_{25}^3$. % {https://www.deltami.edu.pl/2025/03/wycinanki-i-ukladanki/} % lang-pl
di Miłosz Kwiatkowski in $\Delta_{25}^3$. % lang-it

\begin{definition}
    Dwa wielokąty nazywamy równoważnymi przez doklejanie, jeśli możemy dołączyć do nich pary przystających kawałków tak, by końcowe figury były przystające przez dodawanie. % lang-pl
    Due poligoni si dicono equicompletabili se a ciascuno di essi possiamo aggiungere coppie di pezzi congruenti in modo che le figure finali risultino equiscomponibili. % lang-it
    Formalnie: istnieją przystające przez dodawanie wielokąty $R \cong_+ S$ oraz pewne wielokąty $U_1, U_2, \ldots, U_n$ i $V_1, V_2, \ldots, V_n$ takie, że: $U_1 \cong V_i$, $R = P \sqcup U_1 \sqcup \ldots \sqcup U_n$ oraz $S = Q \sqcup V_1 \sqcup \ldots \sqcup V_n$. % lang-pl
    Formalmente: esistono poligoni equiscomponibili $R \cong_+ S$ e alcuni poligoni $U_1, U_2, \ldots, U_n$ e $V_1, V_2, \ldots, V_n$ tali che: $U_1 \cong V_i$, $R = P \sqcup U_1 \sqcup \ldots \sqcup U_n$ e $S = Q \sqcup V_1 \sqcup \ldots \sqcup V_n$. % lang-it
    Oznaczamy to $P \cong_- Q$. % lang-pl
    Lo denotiamo con $P \cong_- Q$. % lang-it
\end{definition}

\begin{proposition}
    \label{congruent_addition_subtraction}%
    Prawdziwe są implikacje: % lang-pl
    Valgono le seguenti implicazioni: % lang-it
    \begin{itemize}
        \item jeśli $P \cong Q$, to $P \cong_+ Q$. % lang-pl
        \item se $P \cong Q$, allora $P \cong_+ Q$. % lang-it
        \item jeśli $P \cong_+ Q$, to $P \cong_- Q$. % lang-pl
        \item se $P \cong_+ Q$, allora $P \cong_- Q$. % lang-it
        \item jeśli $P \cong_- Q$, to $P$ oraz $Q$ są równopolowe. % lang-pl
        \item se $P \cong_- Q$, allora $P$ e $Q$ sono equivalenti. % lang-it
    \end{itemize}
\end{proposition}

Konieczność dowodu tego twierdzenia dostrzeże W. H. Jackson w 1912 roku, szczegóły można znaleźć u Evesa \cite[s. 198, 199]{eves1_1972} (gdzie pokazuje też, że relacje $\cong_+$ i $\cong_-$ są przechodnie). % lang-pl
La necessità di dimostrare questo fatto sarà notata da W. H. Jackson nel 1912; i dettagli si possono trovare in Eves \cite[p. 198, 199]{eves1_1972} (dove si mostra anche che le relazioni $\cong_+$ e $\cong_-$ sono transitive). % lang-it

Eves wprowadza pojęcie wierzchołka rzutującego (z angielskiego \emph{projecting vertex}, czyli taki, który można oddzielić prostą od pozostałych wierzchołków) i pokazuje, że każdy wielokąt ma co najmniej trzy niewspółliniowe wierzchołki rzutujące. % lang-pl
Eves introduce il concetto di vertice proiettante (dall'inglese \emph{projecting vertex}, cioè un vertice che si può separare dagli altri vertici mediante una retta) e mostra che ogni poligono ha almeno tre vertici proiettanti non allineati. % lang-it
Jest to potrzebne do uzasadnienia, że: % lang-pl
Ciò è necessario per giustificare che: % lang-it

\begin{proposition}
    Każdy wielokąt posiada triangulację: rozcięcie na trójkąty, których wierzchołki są wierzchołkami wyjściowego wielokąta. % lang-pl
    Ogni poligono possiede una triangolazione: una dissezione in triangoli i cui vertici sono vertici del poligono di partenza. % lang-it
\end{proposition}

W przypadku wielokątów wypukłych możemy żądać nawet, by wszystkie trójkąty miały wspólny wierzchołek. % lang-pl
Nel caso di poligoni convessi possiamo persino richiedere che tutti i triangoli abbiano un vertice comune. % lang-it

\begin{proposition}
    Każdy trójkąt jest równoważny przez rozcinanie z równopolowym prostokątem o boku takim samym jak najdłuższy bok trójkąta. % lang-pl
    Ogni triangolo è equiscomponibile con un rettangolo di area uguale avente un lato uguale al lato più lungo del triangolo. % lang-it
\end{proposition}

\begin{corollary}
    Każdy prostokąt jest równoważny przez rozcinanie z kwadratem o takim samym polu. % lang-pl
    Ogni rettangolo è equiscomponibile con un quadrato della stessa area. % lang-it
\end{corollary}

\begin{corollary}
    Każdy wielokąt jest równoważny przez rozcinanie z kwadratem o takim samym polu. % lang-pl
    Ogni poligono è equiscomponibile con un quadrato della stessa area. % lang-it
\end{corollary}

Doszliśmy wreszcie do wyniku Williama Wallace'a (z 1807 roku), znalezionego niezależnie przez Farkasa Bolyaia (1833) i Paula Gerwiena (1835): % lang-pl
Siamo infine giunti al risultato di William Wallace (del 1807), trovato indipendentemente da Farkas Bolyai (1833) e Paul Gerwien (1835): % lang-it
% P. Gerwien, Zerschneidung jeder beliebigen Anzahl von gleichen geradlinigen Figuren in dieselben Stu¨ cke, Journal fu¨ r die reine und angewandte Mathematik 10(1833), 228–234.

% W. Bolyai de Bolya, Tentamen. Iuventutem studiosam in ele- menta matheseos purae, elementaris ac sublı`mioris, methodo intuitiva, evidentiaque huic propria, introducendi. Cum appendice triplici, Maros Va´ sa´ rhelyini; tomus primus 1932, tomus secudus 1933.

\begin{theorem}
[Wallace'a-Bolyaia-Gerwiena] % lang-pl
[di Wallace-Bolyai-Gerwien] % lang-it
\label{twierdzenie_wallace_bolyai_gerwien}%
    Następujące warunki są równoważne: % lang-pl
    Le seguenti condizioni sono equivalenti: % lang-it
    \begin{itemize}
    \item dwie figury są równopolowe, % lang-pl
    \item due figure sono equivalenti, % lang-it
    \item dwie figury są równoważne przez rozcinanie. % lang-pl
    \item due figure sono equiscomponibili. % lang-it
    \end{itemize}
    Punkty rozcięcia można wyznaczyć przy pomocy samego cyrkla i linijki. % lang-pl
    I punti di taglio si possono determinare usando solo compasso e riga. % lang-it
\end{theorem}
% https://en.wikipedia.org/wiki/Wallace-Bolyai-Gerwien_theorem

Laczkovich wzmocni później ten wynik: każdy wielokąt można rozciąć na kawałki i złożyć z nich równopolowy kwadrat tak, by kawałki były przesuwane, ale nie obracane. % lang-pl
Laczkovich rafforzerà in seguito questo risultato: ogni poligono può essere tagliato in pezzi e ricomposto in un quadrato equivalente in modo tale che i pezzi vengano traslati, ma non ruotati. % lang-it
% M. Laczkovich, Equidecomposability and discrepancy: a solution to Tarski’s circle squaring problem, Journal fu¨ r die reine und angewandte Mathematik 404(1990), 77–117.
% M. Laczkovich, Paradoxical decompositions: a survey of recent results, Proc. First European Congress of Mathematics, Vol. II (Paris, 1992), Progress in Mathematics 120, Birkha¨ user, Basel 1994

\subsection{Miara Lebesgue'a} % lang-pl
\subsection{La misura di Lebesgue} % lang-it

Skończymy na odległych czasach. % lang-pl
Concluderemo con tempi lontani. % lang-it
Tysiące lat po Euklidesie pojęcie rozmiaru (nie tylko długości, ale przede wszystkim pola figur płaskich oraz objętości brył) sformalizuje się przez miarę Jordana, albo jej uogólnienie --- miarę Lebesgue'a. % lang-pl
Migliaia di anni dopo Euclide, il concetto di grandezza (non solo la lunghezza, ma soprattutto l'area delle figure piane e il volume dei solidi) si formalizzerà attraverso la misura di Jordan, oppure la sua generalizzazione: la misura di Lebesgue. % lang-it
O mierze wspomina Hartshorne \cite[s. 205]{hartshorne_2000}; jest to funkcja, która przyporządkowuje każdej figurze $P$ pewną nieujemną liczbę $[P]$ (zwaną polem) taką, że: % lang-pl
Della misura parla Hartshorne \cite[p. 205]{hartshorne_2000}; si tratta di una funzione che associa a ogni figura $P$ un numero non negativo $[P]$ (detto area) tale che: % lang-it
\begin{itemize}
    \item pole każdego trójkąta jest dodatnie: $[\triangle] > 0$; % lang-pl
    \item l'area di ogni triangolo è positiva: $[\triangle] > 0$; % lang-it
    \item przystające trójkąty mają równe pola; % lang-pl
    \item triangoli congruenti hanno aree uguali; % lang-it
    \item jeśli dwie figury $P, Q$ nie nachodzą na siebie, to $[P \cup Q] = [P] + [Q]$. % lang-pl
    \item se due figure $P, Q$ non si sovrappongono, allora $[P \cup Q] = [P] + [Q]$. % lang-it
\end{itemize}

Z istnienia takiej funkcji wynika aksjomat (Z), wysłowiony przez nas niżej. % lang-pl
Dall'esistenza di una tale funzione segue l'assioma (Z), da noi enunciato più sotto. % lang-it
Guzicki \cite[s. 38-68]{guzicki_2021} również wspomni o mierze Jordana. % lang-pl
Anche Guzicki \cite[p. 38-68]{guzicki_2021} menzionerà la misura di Jordan. % lang-it
Zauważy jednak, że w polskich szkołach każdą figurę można rozciąć na wielokąty i wycinki kołowe, co prowadzi do prostszej aksjoamtyzacji (i naturalnie do dowodu twierdzenia \ref{twierdzenie_wallace_bolyai_gerwien}, Wallace'a-Bolyaia-Gerwiena). % lang-pl
Osserverà però che, nelle scuole polacche, ogni figura può essere ritagliata in poligoni e settori circolari, il che porta a un'assiomatizzazione più semplice (e naturalmente alla dimostrazione del teorema \ref{twierdzenie_wallace_bolyai_gerwien}, di Wallace-Bolyai-Gerwien). % lang-it

Pole jest wyznaczone jednoznacznie z dokładnością do stałej. % lang-pl
L'area è determinata in modo univoco a meno di una costante. % lang-it
Wartość stałej ustala się zazwyczaj przez przyjęcie, że trójkąt o podstawie $a$ oraz wysokości $h$ ma pole $\frac 1 2 ah$. % lang-pl
Il valore della costante si fissa di solito assumendo che un triangolo di base $a$ e altezza $h$ abbia area $\frac 1 2 ah$. % lang-it

\subsection{Zasada Cavalieriego} % lang-pl
\subsection{Il principio di Cavalieri} % lang-it

Sekcję zamkniemy ciekawostką związaną z zasadą Cavalieriego. % lang-pl
Chiuderemo la sezione con una curiosità legata al principio di Cavalieri. % lang-it
Eves pokaże w 1991 roku, że każde dwa równopolowe trójkąty można umieścić na płaszczyźnie tak, by każda prosta równoległa do danej przecinała obydwa na odcinku o tej samej długości, wspomni o tym Jarosław Górnicki w $\Delta_{11}^{11}$. % lang-pl
Eves mostrerà nel 1991 che due triangoli equivalenti qualsiasi si possono sempre collocare nel piano in modo che ogni retta parallela a una retta data intersechi entrambi in un segmento della stessa lunghezza; ne parlerà Jarosław Górnicki in $\Delta_{11}^{11}$. % lang-it
% {https://www.deltami.edu.pl/2011/11/osobliwosc-trojkatow/}

%
